\chapter{Experimental Methodology}
\label{ch:methodology}

The evaluation's credibility rests on rules that were fixed before results
existed. This chapter states them; Chapter~\ref{ch:evaluation} reports what
happened under them --- including the hypotheses that failed.

\section{Pre-registration}
\label{sec:method-prereg}

Every confirmatory campaign's protocol --- hypotheses, sole treatment,
analysis script, and stopping rule --- was committed and pushed to the
public repository \emph{before its first run}. The first seven were
\evfile{PREREG\_V2.md}, \evfile{PREREG\_V3.md}, \evfile{PREREG\_TRACE.md},
\evfile{PREREG\_TRACE2.md}, \evfile{PREREG\_VTC.md}, plus the v1 gate and
the Phase-4 acceptance frozen in the roadmap documents; forty-eight
\evfile{PREREG\_*.md} files exist at the time of writing, each with its
push timestamp as the anchor, all of them mirrored --- frozen, public and
DOI-bearing --- on the Open Science Framework
(\url{https://doi.org/10.17605/OSF.IO/DYZKV}; index \evfile{OSF\_REGISTRATION.md}),
a mirror made after the campaigns ran and stated as such, and a gate
(\evfile{scripts/check\_preregs.py}) proves from the git history that no
registered text was altered after its result existed. Protocol changes
forced by feasibility (e.g.\ the LMSYS similarity matrix that would not fit
in 40\,TB) were committed as \emph{declared pre-run amendments} before the
gated data was downloaded --- the amendment chain is the pre-registration's
honest continuation, visible in git history. Closed campaigns are
immutable: stopping rules forbid re-running, widening, or post-hoc tuning,
and campaign result files are never edited after their run. When a
confound is found \emph{after} a campaign closes --- as three were, in an
audit of the baseline implementations (Section~\ref{sec:eval-adjudication})
--- the response is a new pre-registration whose outcome for the published
claim is committed before the re-scoring runs, and a new record beside the
old one; the published record stands, and the thesis quotes the surviving
claim.

\section{Baselines, tuned on our own objective}
\label{sec:method-baselines}

Five baselines: HPA-style reactive scaling, KEDA-style event-driven scaling,
a FIRM-style replica controller~\citep{firm2020}, GPTCache-style fixed
semantic caching with LRU~\citep{gptcache2023}, and static provisioning;
later campaigns add a VTC-style least-service-first replica divider
\citep{vtc2024} and iso-cost variants. Every baseline is grid-search tuned
\emph{on the composite objective this thesis is graded on}, with sweeps
committed (\evfile{eval/baselines/TUNING.md}, grids in
\evfile{eval/baselines/grids/}). Grids are bounded at a vendor-sane envelope
because $J$ improves monotonically toward over-provisioning --- an unbounded
grid would tune every baseline into a cost blowout and flatter \polyforge{}.
Iso-cost variants exist to remove the ``it just spends less'' objection, and
the gate they impose is reported FAIL-honest where it fails.

\section{Substrates and separation rules}
\label{sec:method-substrates}

Three substrates carry the evidence, and they are never mixed in one table:

\begin{enumerate}
  \item \textbf{The simulator} (\evfile{research/jcac\_sim}): discrete-event
        decision-quality substrate. Its congestion form was calibrated
        against a real inference server --- the fitted exponent
        ($a{=}0.86$ vs.\ assumed $1.0$) means the model over-penalizes high
        utilization, i.e.\ errs \emph{against} \polyforge{}'s lean postures
        --- and its tier constants were checked against a real GPU bench
        (measured latency ratios 1:1.5:16.6; same ordering as the priced
        1:10:100, compressed scale, disclosed).
  \item \textbf{Real data replayed through the simulator}: BurstGPT
        v2.0~\citep{burstgpt2024} demand (10{,}632{,}194 requests, 335
        days), Azure LLM inference 2024~\citep{azurellmtrace2024} demand
        (44.1M requests), and LMSYS-Chat-1M prompts~\citep{lmsyschat1m2024}
        under a real MiniLM encoder~\citep{minilm2020,sbert2019}.
  \item \textbf{A live kind cluster}: real scheduler, real HPA, real load.
        Claims from this substrate are \emph{ordinal only} --- does the
        simulator's ranking reproduce? --- never absolute comparisons.
\end{enumerate}

Replay dollars are model-scale; the claim is always the ranking. Samples
from different replays are never pooled ($n{=}16$ and $n{=}96$ are separate
experiments with separate pre-registrations).

\section{Statistics}
\label{sec:method-stats}

Runs are paired by seed: every system sees identical demand realizations,
including arrival jitter drawn once per window and shared across systems.
Significance is Wilcoxon/paired tests at the pre-registered gate; but at
hundreds of paired seeded simulator runs, vanishing p-values measure the
simulator's determinism, not the effect. The \emph{citable unit} is
therefore the paired effect size $\dz$ with a 95\% bootstrap confidence
interval (\evfile{EFFECT\_SIZES.md} restates the real-trace campaigns this
way; all $n{=}96$ confirmatory intervals exclude zero, and the $+0.07$
violation trade carries its own CI). P-values appear once, as prereg gate
outcomes, and are never headlined below ${\sim}10^{-6}$. A
declared-before-run weight-sensitivity sweep perturbs the objective weights
across 425 cells spanning all campaigns: the win direction never flips, and
416/425 cells keep $p<0.01$ (\evfile{SENSITIVITY\_J.md}).

\section{Metrics}
\label{sec:method-metrics}

Cost (USD: replica-hours at \$0.048/replica-hr; model-tier calls at market
API ratios 1:10:100), SLO violation rate, semantic-cache hit rate, Jain's
index over per-tenant SLO satisfaction, and the composite
\begin{equation}
J = \text{cost} + 2\cdot\text{violation} + 0.5\cdot(1-\text{Jain}),
\end{equation}
normalized per tenant-step --- the objective every baseline was tuned on.
Latency is reported at \textbf{p95, not p99}: the simulator's latency
estimator is a queueing approximation whose tail beyond p95 is not credible,
and pretending otherwise would be false precision. This is a documented
deviation; the live path measures wall-time latency and can report any
percentile.

\section{The workload matrix}
\label{sec:method-matrix}

The headline matrix is a blocked factorial: 5 workload classes (from the
taxonomy) $\times$ 4 tenant-mix profiles $\times$ 3 cluster sizes $\times$ 6
systems $\times$ 5 repetitions $=$ 1{,}800 runs, plus 500 ablation runs.
Eight tenants per cluster is the studied factor (composition, not
population); the planner's tenant-count scaling is measured separately
(Section~\ref{sec:impl-scaling}). The v3 overload cells were sized by
\emph{arithmetic on committed model constants} --- burst climbs exceeding
the shared $\pm2$-replicas-per-interval actuation clamp --- pinned by unit
tests, not tuned by trial runs, with a same-envelope
\texttt{ramp\_gentle} control cell where the hypothesized advantage should
not (and does not) appear.

\section{Honest-nulls policy}
\label{sec:method-nulls}

Negative and failed results are published with the same prominence as wins,
in a dedicated ledger (\evfile{RESULTS\_MASTER.md}): the two confirmatory
SLO failures (v2 H1, v3 H1$'$), the underpowered first replay (HT), the
fairness-term null ($\gamma$), the seasonal-forecasting non-transfer from
synthetic to real demand, the structurally impossible raw per-metric gate,
and the cache hit-quality result that most $\tau{=}0.85$ hits fail the
response-agreement proxy. Goalpost-moving is invisible when it happens in
private; this project's position is that it happened in public, in the git
history, where an examiner can diff it.
